\begin{thaiabstract}
ในปี\hspace{0.15cm}ค.ศ.\hspace{0.1cm}1976\hspace{0.2cm}ดอว์สันและโฮวาร์ดได้แสดงว่าเซตของฟังก์ชันหนึ่งต่อหนึ่งทั่วถึงบนเซต $A$ เขียนแทนด้วย $S(A)$ มีขนาดเท่ากับเซตกำลังของ $A$ เขียนแทนด้วย $\mathcal{P}(A)$ เมื่อ $A$ เป็นเซตอนันต์ ในทฤษฎีเซตแซร์เมโล-แฟรงเคล (ZF) ที่มีสัจพจน์การเลือก (AC) แต่หากปราศจาก AC แล้ว จะไม่สามารถสรุปความสัมพันธ์ระหว่างขนาดของ $S(A)$ และ $\mathcal{P}(A)$ ได้สำหรับเซตอนันต์ $A$ ใด ๆ ต่อมาได้มีการศึกษาความสัมพันธ์ของ $|S(A)|$ กับจำนวนเชิงการนับอื่น ๆ เช่นกัน อาทิ ได้มีการแสดงว่า ``$|\seq(A)|\neq|S(A)|\neq|\seq^{1-1}(A)|$ สำหรับเซตอนันต์ $A$ ใด ๆ" เป็นผลลัพธ์ที่ดีที่สุดใน ZF ที่ปราศจาก AC โดยที่ $\seq(A)$ และ $\seq^{1-1}(A)$ คือเซตของลำดับจำกัดและเซตของลำดับจำกัดที่ไม่มีการซ้ำของสมาชิกในเซต $A$ ตามลำดับ  ในวิทยานิพนธ์นี้ เราศึกษาเซตของฟังก์ชันหนึ่งต่อหนึ่งและเซตของฟังก์ชันทั่วถึงบนเซต $A$ เขียนแทนด้วย  $I(A)$ และ $J(A)$ ตามลำดับ เราได้แสดงว่า อสมการข้างต้นเป็นจริงเช่นกันเมื่อ $S(A)$ ถูกแทนด้วย $I(A)$ และ $J(A)$  นอกจากนี้ เราจะไม่มีทางพิสูจน์ใน ZF ได้ว่า ``$|I(A)|=|J(A)|$ สำหรับเซตอนันต์ $A$ ใด ๆ"

\end{thaiabstract}